\documentclass[10pt]{article}
\usepackage[T1]{fontenc}
\usepackage[utf8]{inputenc}
\usepackage[french]{babel}
\usepackage{amsmath,amssymb,mathrsfs,mathpazo}
\usepackage{geometry}
\geometry{a4paper,margin=2cm}
\usepackage{graphicx,array,tabularx,tabularray,colortbl,multirow,multicol,booktabs,enumitem,pifont,ulem,fancybox,diagbox,cancel}
\usepackage{pstricks,pstricks-add,pst-plot,pst-tree,pst-node,pst-text,pst-func,pst-math,pst-bspline,pst-3dplot,pst-eucl,pst-solides3d}
\let\pscalculate\undefined
\usepackage{pst-bezier}
\usepackage{tikz,pgfplots,tkz-tab}
\usetikzlibrary{arrows,arrows.meta,patterns,calc,positioning,trees,decorations.markings,shapes.geometric,shapes.misc,shapes.arrows,intersections,decorations.pathmorphing,decorations.pathreplacing,backgrounds,shadows,babel}
\pgfplotsset{compat=1.18}
\usepgfplotslibrary{fillbetween}
\usepackage[np]{numprint}
\usepackage{listings,makecell,pas-tableur,twoopt}
\let\barre\undefined
\usepackage[pstricks]{tablvar}
\definecolor{darkgreen}{RGB}{0,100,0}
\providecommand{\red}{\color{red}}
\providecommand{\blue}{\color{blue}}
\providecommand{\green}{\color{green}}
\providecommand{\black}{\color{black}}
\providecommand{\white}{\color{white}}
\providecommand{\gray}{\color{gray}}
\providecommand{\cyan}{\color{cyan}}
\providecommand{\magenta}{\color{magenta}}
\providecommand{\yellow}{\color{yellow}}
\providecommand{\violet}{\color{violet}}
\providecommand{\orange}{\color{orange}}
\providecommand{\euro}{\texteuro}
\providecommand{\EUR}{\texteuro}
\providecommand{\ieme}{\textsuperscript{e}}
\providecommand{\degres}{\ensuremath{^\circ}}
\providecommand{\degree}{\ensuremath{^\circ}}
\providecommand{\vv}[1]{\overrightarrow{#1}}
\providecommand{\Coord}[1]{#1}
\providecommand{\decosix}{}
\providecommand{\starredbullet}{\ensuremath{\bullet}}
\pagestyle{empty}
\setlength{\parindent}{0pt}
\usepackage{adjustbox}\DeclareUnicodeCharacter{2236}{:}\DeclareUnicodeCharacter{2190}{\ensuremath{\leftarrow}}
\providecommand{\R}{\mathbb{R}}
\providecommand{\N}{\mathbb{N}}
\providecommand{\Z}{\mathbb{Z}}
\providecommand{\Q}{\mathbb{Q}}
\providecommand{\D}{\mathbb{D}}
\providecommand{\C}{\mathbb{C}}
\providecommand{\e}{{\mathrm{e}}}
\providecommand{\vect}[1]{\overrightarrow{#1}}
\providecommand{\vectt}[1]{\overrightarrow{\mathrm{#1}}}
\providecommand{\vvt}[1]{\overrightarrow{\mathrm{#1}}}
\providecommand{\np}[2][]{{#2~#1}}
\providecommand{\vv}[1]{{\overrightarrow{#1}}}
\renewcommand{\d}{{\mathrm{d}}}
\providecommand{\degres}{{^\circ}}
\providecommand{\degre}{{^\circ}}
\providecommand{\blue}{\color{blue}}
\providecommand{\red}{\color{red}}
\providecommand{\enskip}{\quad}
\providecommand{\llbracket}{\lbrack\!\lbrack}
\providecommand{\rrbracket}{\rbrack\!\rbrack}
\providecommand{\euro}{{\text{€}}}
\providecommand{\no}{{\text{n°}}}
\providecommand{\interval}[2]{{[#1;#2]}}
\providecommand{\intervFO}[2]{{[#1;#2[}}
\providecommand{\dotfill}{\ldots}

\providecommand{\ds}{\displaystyle}
\providecommand{\textding}[1]{#1}
\providecommand{\Oij}{\(\mathrm{O};\vec{\imath},\vec{\jmath}\)}
\providecommand{\Oijk}{\(\mathrm{O};\vec{\imath},\vec{\jmath},\vec{k}\)}
\providecommand{\Ouv}{\(\mathrm{O};\vec{u},\vec{v}\)}
\providecommand{\intervFF}[2]{\ensuremath{[#1;#2]}}
\renewcommand{\euro}{\ifmmode\text{\texteuro}\else\texteuro\fi}
\DeclareUnicodeCharacter{00B0}{\ifmmode^{\circ}\else\textdegree\fi}
\DeclareUnicodeCharacter{20AC}{\ifmmode\text{\texteuro}\else\texteuro\fi}
% Énoncés : Éducation nationale. Transcription et corrigés : APMEP, Denis Vergès et contributeurs. Découpage et classement : CoopMaths, Sébastien Lozano, Éric Elter et contributeurs. Adaptation de la présentation : Maths Explorer.
% CC BY-SA : https://coopmaths.fr/www/about/
% https://coopmaths.fr/alea/static/e3c/2020/tex/e3c_2020_00_sujet29_1_cor.tex
\begin{document}
\noindent\textit{APMEP / CoopMaths. CC BY-SA. Présentation adaptée par Maths Explorer.}\par\medskip
\subsection*{1.}

L'arbre pondéré ci-dessous représente une situation où \(A\), \(B\), \(C\) et \(D\) sont des événements d'une expérience aléatoire :

\begin{center}
\begin{tikzpicture}[thick, scale=1.5]
\node (P_-1_0) at (-2,-2.5) {$\phantom{A}$};
\node (P_0_0) at (0,-0.5) {$A$};
\draw (P_-1_0) -- (P_0_0) node[midway, above,sloped] {$0,12$};
\node (P_1_0) at (2,-0) {$D$};
\draw (P_0_0) -- (P_1_0) node[midway, above] {$0,5$};
\node (P_1_1) at (2,-1) {$\overline{D}$};
\draw (P_0_0) -- (P_1_1) node[midway, below] {$0,5$};
\node (P_0_2) at (0,-2.5) {$B$};
\draw (P_-1_0) -- (P_0_2) node[midway, above] {$0,24$};
\node (P_1_2) at (2,-2) {$D$};
\draw (P_0_2) -- (P_1_2) node[midway, above] {$0,2$};
\node (P_1_3) at (2,-3) {$\overline{D}$};
\draw (P_0_2) -- (P_1_3) node[midway, below] {$0,8$};
\node (P_0_4) at (0,-4.5) {$C$};
\draw (P_-1_0) -- (P_0_4) node[midway, below,sloped] {$0,64$};
\node (P_1_4) at (2,-4) {$D$};
\draw (P_0_4) -- (P_1_4) node[midway, above] {$0,1$};
\node (P_1_5) at (2,-5) {$\overline{D}$};
\draw (P_0_4) -- (P_1_5) node[midway, below] {$0,9$};
\end{tikzpicture}
\end{center}

D'après la loi des probabilités totales :

\(
p(D) = p(A \cap D) + p(B \cap D) + p(C \cap D),
\)

\(
p(D) = p(A) \times p_A(D) + p(B) \times p_B(D) + p(C) \times p_C(D),
\)

\(
p(D) = 0,12 \times 0,5 + 0,24 \times 0,2 + 0,64 \times 0,1 = 0,06 + 0,048 + 0,064 = 0,172.
\)

\subsection*{2.}

Pour le trinôme \(-2x^2 - 5x + 3\), on a :
\(
\Delta = 25 + 4 \times 2 \times 3 = 25 + 24 = 49 = 7^2 > 0.
\)

Ce trinôme a deux racines :
\(
x_1 = \dfrac{5 + 7}{-4} = -3 \quad \text{et} \quad x_2 = \dfrac{5 - 7}{-4} = \dfrac{1}{2}.
\)

On sait que ce trinôme est négatif sauf sur l'intervalle \(\left] -3 \,;\, \dfrac{1}{2} \right[\).

Donc :
\(
S = \left] -\infty \,;\, -3 \right[ \cup \left] \dfrac{1}{2} \,;\, +\infty \right[.
\)

\subsection*{3.}

Un vecteur directeur de la droite a pour coordonnées \(\begin{pmatrix} 8 \\ 2 \end{pmatrix}\),
donc un vecteur normal a pour coordonnées \(\begin{pmatrix} -2 \\ 8 \end{pmatrix}\) ou \(\begin{pmatrix} 1 \\ -4 \end{pmatrix}\).

\subsection*{4.}

\begin{align*}
&M(x\,;\,y) \in \mathcal{C}(A\,;\,R = 2)\\
\iff& (AM^2 = 2^2\\
\iff &(x - (-2))^2 + (y - (-4))^2 = 4\\
\iff &(x + 2)^2 + (y + 4)^2 = 4\\
\iff &x^2 + 4 + 4x + y^2 + 16 + 8y = 4\\
\iff &x^2 + y^2 + 4x + 8y + 16 = 0
\end{align*}

\subsection*{5.}

On a effectivement :
\[
u_1 = u_0 + 2 \times 0 - 3 = 1 + 0 - 3 = -2,
\]
\[
u_2 = u_1 + 2 \times 1 - 3 = -2 + 2 - 3 = -3.
\]

\end{document}
